W niniejszym rozdziale przedstawiono analizę wyników testów skalowalności przeprowadzonych dla
trzech implementacji językowych: Python, JavaScript oraz Java. Badania objęły dwie wersje
architektoniczne aplikacji: monolityczną oraz mikroserwisową wykorzystującą \texttt{gRPC} jako
technologię komunikacji międzyserwisowej. Analizie poddano dwa mechanizmy zwiększania zasobów:
skalowanie wertykalne, polegające na zwiększaniu parametrów pojedynczej instancji maszyny,
oraz skalowanie horyzontalne, polegające na zwiększaniu liczby instancji aplikacji lub replik
mikroserwisów.

Wyniki zestawiono zarówno w postaci wykresów przedstawiających maksymalną osiąganą przepustowość
w funkcji narastającego obciążenia, jak i w tabelach zbiorczych prezentujących maksymalne wartości
przepustowości oraz przyrosty wydajności. Celem analizy było określenie, który mechanizm skalowania
jest skuteczniejszy dla poszczególnych języków i architektur oraz jak duży wpływ na końcowy rezultat
ma sama technologia wykonania aplikacji.

\section{Wyniki skalowania wertykalnego}

W testach skalowania wertykalnego porównano trzy konfiguracje zasobów obliczeniowych:
\texttt{std-2}, \texttt{std-4} oraz \texttt{std-8}. Oznacza to odpowiednio bazowy poziom zasobów,
konfigurację dwukrotnie większą oraz konfigurację czterokrotnie większą względem punktu odniesienia.
Tabela~\ref{tab:vertical-scaling-results} przedstawia zbiorcze wyniki dla wszystkich badanych
implementacji.

\begin{table}[htbp]
\centering
\caption{Wyniki skalowania wertykalnego dla badanych implementacji.}
\label{tab:vertical-scaling-results}
\renewcommand{\arraystretch}{1.15}
\resizebox{\textwidth}{!}{%
\begin{tabular}{|c|l|l|l|r|r|r|r|r|r|}
\hline
\textbf{Lp.} & \textbf{Język} & \textbf{Architektura} & \textbf{Komunikacja} & \textbf{std-2} & \textbf{std-4} & \textbf{std-8} & \textbf{std-4/std-2} & \textbf{std-8/std-2} & \textbf{std-8/std-4} \\
 &  &  &  & \textbf{[req/s]} & \textbf{[req/s]} & \textbf{[req/s]} & \textbf{[\%]} & \textbf{[\%]} & \textbf{[\%]} \\
\hline
1 & Python & monolityczna & - & 77,2 & 145,6 & 216,4 & 88,6\% & 180,3\% & 48,6\% \\
2 & JavaScript & monolityczna & - & 1058,8 & 1096,8 & 1173,4 & 3,6\% & 10,8\% & 7,0\% \\
3 & Java & monolityczna & - & 1841,6 & 3630,6 & 7275 & 97,1\% & 295,0\% & 100,4\% \\
4 & Python & mikroserwisowa & gRPC & 45,8 & 55,8 & 63,6 & 21,8\% & 38,9\% & 14,0\% \\
5 & JavaScript & mikroserwisowa & gRPC & 530,4 & 699,2 & 946,4 & 31,8\% & 78,4\% & 35,4\% \\
6 & Java & mikroserwisowa & gRPC & 1354 & 2523,6 & 4835,8 & 86,4\% & 257,1\% & 91,6\% \\
\hline
\end{tabular}%
}
\end{table}

\begin{figure}[htbp]
\centering
\begin{minipage}{0.49\textwidth}
    \centering
    \includegraphics[width=\textwidth]{images/scaling/vertical-scaling/python-monolith-images_vertical_scaling}

    \small (a) Python -- architektura monolityczna
\end{minipage}
\hfill
\begin{minipage}{0.49\textwidth}
    \centering
    \includegraphics[width=\textwidth]{images/scaling/vertical-scaling/python-microservices-grpc-images_vertical_scaling}

    \small (b) Python -- architektura mikroserwisowa z \texttt{gRPC}
\end{minipage}
\caption{Wyniki skalowania wertykalnego dla implementacji w języku Python.}
\label{fig:vertical-python}
\end{figure}

\begin{figure}[htbp]
\centering
\begin{minipage}{0.49\textwidth}
    \centering
    \includegraphics[width=\textwidth]{images/scaling/vertical-scaling/javascript-monolith-images_vertical_scaling}

    \small (a) JavaScript -- architektura monolityczna
\end{minipage}
\hfill
\begin{minipage}{0.49\textwidth}
    \centering
    \includegraphics[width=\textwidth]{images/scaling/vertical-scaling/javascript-microservices-grpc-images_vertical_scaling}

    \small (b) JavaScript -- architektura mikroserwisowa z \texttt{gRPC}
\end{minipage}
\caption{Wyniki skalowania wertykalnego dla implementacji w języku JavaScript.}
\label{fig:vertical-javascript}
\end{figure}

\begin{figure}[htbp]
\centering
\begin{minipage}{0.49\textwidth}
    \centering
    \includegraphics[width=\textwidth]{images/scaling/vertical-scaling/java-monolith-images_vertical_scaling}

    \small (a) Java -- architektura monolityczna
\end{minipage}
\hfill
\begin{minipage}{0.49\textwidth}
    \centering
    \includegraphics[width=\textwidth]{images/scaling/vertical-scaling/java-microservices-grpc-images_vertical_scaling}

    \small (b) Java -- architektura mikroserwisowa z \texttt{gRPC}
\end{minipage}
\caption{Wyniki skalowania wertykalnego dla implementacji w języku Java.}
\label{fig:vertical-java}
\end{figure}

Analiza wyników skalowania wertykalnego pokazuje wyraźnie, że największe korzyści z dokładania
zasobów uzyskano dla implementacji w języku Java. W przypadku monolitu przepustowość wzrosła
z poziomu 1841,6~req/s dla konfiguracji \texttt{std-2} do 7275~req/s dla konfiguracji \texttt{std-8},
co oznacza przyrost o 295,0\%. Bardzo podobny trend zaobserwowano dla mikroserwisów w Javie,
gdzie wzrost wyniósł 257,1\%. Oznacza to, że implementacja java efektywnie wykorzystuje dodatkowe
rdzenie procesora i pamięć operacyjną, a skalowanie pionowe przekłada się bezpośrednio na wzrost
przepustowości.

Odmienny obraz widoczny jest dla JavaScriptu. W architekturze monolitycznej wzrost wydajności
po zwiększeniu zasobów był minimalny: z 1058,8~req/s do 1173,4~req/s, czyli jedynie o 10,8\%.
Sugeruje to, że w tej wersji aplikacji głównym ograniczeniem nie był sam przydział CPU lub RAM,
lecz charakter wykonania pojedynczego procesu. Dla mikroserwisów w JavaScript przyrost był
wyraźnie większy i osiągnął 78,4\%, jednak nadal pozostawał istotnie niższy niż w przypadku Javy.

Python osiągnął najniższe wartości bezwzględne przepustowości, ale sam kierunek zmian był
zróżnicowany zależnie od architektury. Monolit zwiększył przepustowość z 77,2~req/s do
216,4~req/s, czyli o 180,3\%, natomiast wersja mikroserwisowa jedynie z 45,8~req/s do 63,6~req/s,
co daje wzrost rzędu 38,9\%. Oznacza to, że w Pythonie skalowanie pionowe poprawia wynik, lecz
w przypadku wersji rozproszonej poprawa ta jest ograniczona.

Z punktu widzenia porównania architektur można zauważyć, że dla każdego języka monolit osiągał
wyższe wartości maksymalnej przepustowości niż odpowiadająca mu implementacja mikroserwisowa
z \texttt{gRPC}. Różnica ta była szczególnie wyraźna dla JavaScriptu oraz Pythona. W przypadku Javy
obie architektury reagowały na skalowanie pionowe bardzo dobrze, jednak monolit nadal zachowywał
przewagę bezwzględną. Wskazuje to, że dodatkowy narzut komunikacji międzyserwisowej pozostaje
istotnym czynnikiem ograniczającym wydajność, nawet wtedy, gdy aplikacja otrzymuje większe
zasoby sprzętowe.

\section{Wyniki skalowania horyzontalnego}

W testach skalowania horyzontalnego analizowano trzy warianty wdrożenia: 1, 2 oraz 4 instancje.
W przypadku monolitu oznaczało to uruchomienie wielu kopii całej aplikacji, natomiast w przypadku
architektury mikroserwisowej zwiększano liczbę replik każdego serwisu. Tabela~\ref{tab:horizontal-scaling-results}
przedstawia zestawienie maksymalnych przepustowości oraz procentowych przyrostów wydajności.

\begin{table}[htbp]
\centering
\caption{Wyniki skalowania horyzontalnego dla badanych implementacji.}
\label{tab:horizontal-scaling-results}
\renewcommand{\arraystretch}{1.15}
\resizebox{\textwidth}{!}{%
\begin{tabular}{|c|l|l|l|r|r|r|r|r|r|}
\hline
\textbf{Lp.} & \textbf{Język} & \textbf{Architektura} & \textbf{Komunikacja} & \textbf{1 inst.} & \textbf{2 inst.} & \textbf{4 inst.} & \textbf{2/1} & \textbf{4/1} & \textbf{4/2} \\
 &  &  &  & \textbf{[req/s]} & \textbf{[req/s]} & \textbf{[req/s]} & \textbf{[\%]} & \textbf{[\%]} & \textbf{[\%]} \\
\hline
1 & Python & monolityczna & - & 79 & 150,4 & 303,2 & 90,4\% & 283,8\% & 101,6\% \\
2 & JavaScript & monolityczna & - & 1017,4 & 1385 & 2318,8 & 36,1\% & 127,9\% & 67,4\% \\
3 & Java & monolityczna & - & 1817,2 & 3390 & 6344,4 & 86,6\% & 249,1\% & 87,2\% \\
4 & Python & mikroserwisowa & gRPC & 46,6 & 58,6 & 79,6 & 25,8\% & 70,8\% & 35,8\% \\
5 & JavaScript & mikroserwisowa & gRPC & 559,2 & 920 & 1925,6 & 64,5\% & 244,3\% & 109,3\% \\
6 & Java & mikroserwisowa & gRPC & 1355,8 & 2292,4 & 5007,8 & 69,1\% & 269,4\% & 118,5\% \\
\hline
\end{tabular}%
}
\end{table}

\begin{figure}[htbp]
\centering
\begin{minipage}{0.49\textwidth}
    \centering
    \includegraphics[width=\textwidth]{images/scaling/horizontal-scaling/python-monolith-images_horizontal_scaling}

    \small (a) Python -- architektura monolityczna
\end{minipage}
\hfill
\begin{minipage}{0.49\textwidth}
    \centering
    \includegraphics[width=\textwidth]{images/scaling/horizontal-scaling/python-microservices-grpc-images_horizontal_scaling}

    \small (b) Python -- architektura mikroserwisowa z \texttt{gRPC}
\end{minipage}
\caption{Wyniki skalowania horyzontalnego dla implementacji w języku Python.}
\label{fig:horizontal-python}
\end{figure}

\begin{figure}[htbp]
\centering
\begin{minipage}{0.49\textwidth}
    \centering
    \includegraphics[width=\textwidth]{images/scaling/horizontal-scaling/javascript-monolith-images_horizontal_scaling}

    \small (a) JavaScript -- architektura monolityczna
\end{minipage}
\hfill
\begin{minipage}{0.49\textwidth}
    \centering
    \includegraphics[width=\textwidth]{images/scaling/horizontal-scaling/javascript-microservices-grpc-images_horizontal_scaling}

    \small (b) JavaScript -- architektura mikroserwisowa z \texttt{gRPC}
\end{minipage}
\caption{Wyniki skalowania horyzontalnego dla implementacji w języku JavaScript.}
\label{fig:horizontal-javascript}
\end{figure}

\begin{figure}[htbp]
\centering
\begin{minipage}{0.49\textwidth}
    \centering
    \includegraphics[width=\textwidth]{images/scaling/horizontal-scaling/java-monolith-images_horizontal_scaling}

    \small (a) Java -- architektura monolityczna
\end{minipage}
\hfill
\begin{minipage}{0.49\textwidth}
    \centering
    \includegraphics[width=\textwidth]{images/scaling/horizontal-scaling/java-microservices-grpc-images_horizontal_scaling}

    \small (b) Java -- architektura mikroserwisowa z \texttt{gRPC}
\end{minipage}
\caption{Wyniki skalowania horyzontalnego dla implementacji w języku Java.}
\label{fig:horizontal-java}
\end{figure}

Wyniki przedstawione w tabeli~\ref{tab:horizontal-scaling-results} pokazują, że skalowanie
horyzontalne okazało się bardzo skuteczne przede wszystkim dla implementacji Java oraz
JavaScript w architekturze mikroserwisowej. Największy procentowy wzrost względem konfiguracji
bazowej osiągnęła mikroserwisowa implementacja w Javie: z 1355,8~req/s do 5007,8~req/s,
co oznacza poprawę o 269,4\%. Bardzo dobry rezultat uzyskał również monolit w Javie,
który zwiększył przepustowość o 249,1\%.

W przypadku JavaScriptu szczególnie istotne jest zachowanie wersji mikroserwisowej.
Przejście z jednej instancji do czterech zwiększyło przepustowość z 559,2~req/s do 1925,6~req/s,
czyli o 244,3\%. Jest to rezultat zdecydowanie lepszy niż w analogicznym teście skalowania
wertykalnego. Taki układ wyników wskazuje, że dla tej implementacji bardziej opłacalne jest
rozmnażanie instancji niż wzmacnianie pojedynczego węzła.

Python pozostaje najsłabszy pod względem wartości bezwzględnych, jednak także tutaj można
zaobserwować poprawę po zwiększaniu liczby instancji. Monolit osiągnął wzrost o 283,8\%,
z 79 do 303,2~req/s, natomiast wersja mikroserwisowa poprawiła wynik o 70,8\%.
Oznacza to, że nawet dla technologii o relatywnie niskiej bazowej przepustowości zwiększanie
liczby instancji może być skuteczną drogą poprawy wydajności.

Warto zwrócić uwagę na fakt, że przyrosty przy przejściu z 2 do 4 instancji nie zawsze są
proporcjonalne do wzrostu liczby replik. Zjawisko to jest naturalne dla systemów rozproszonych,
ponieważ wraz z liczbą instancji rośnie również koszt koordynacji, komunikacji oraz równoważenia
obciążenia. Mimo tego dla większości badanych przypadków skalowanie horyzontalne pozostawało
efektywnym mechanizmem zwiększania przepustowości.

\section{Porównanie skalowania wertykalnego i horyzontalnego}

Aby ocenić, który mechanizm zwiększania zasobów był korzystniejszy, zestawiono ze sobą
wyniki dla konfiguracji odpowiadających dwukrotnemu oraz czterokrotnemu wzrostowi zasobów.
W praktyce porównano wariant 2 instancji z wariantem \texttt{std-4} oraz wariant 4 instancji
z wariantem \texttt{std-8}. Zestawienie przedstawiono w tabeli~\ref{tab:horizontal-vs-vertical-results}
oraz na rysunkach~\ref{fig:compare-python}--\ref{fig:compare-java}.

\begin{table}[htbp]
\centering
\caption{Porównanie skuteczności skalowania horyzontalnego i wertykalnego. Dodatnia wartość
oznacza przewagę skalowania wertykalnego, a ujemna przewagę skalowania horyzontalnego.}
\label{tab:horizontal-vs-vertical-results}
\renewcommand{\arraystretch}{1.15}
\resizebox{\textwidth}{!}{%
\begin{tabular}{|c|l|l|l|r|r|r|r|r|r|}
\hline
\textbf{Lp.} & \textbf{Język} & \textbf{Architektura} & \textbf{Komunikacja} & \textbf{H 2$\times$} & \textbf{V 2$\times$} & \textbf{V/H 2$\times$} & \textbf{H 4$\times$} & \textbf{V 4$\times$} & \textbf{V/H 4$\times$} \\
 &  &  &  & \textbf{[req/s]} & \textbf{[req/s]} & \textbf{[\%]} & \textbf{[req/s]} & \textbf{[req/s]} & \textbf{[\%]} \\
\hline
1 & Python & monolityczna & - & 150,4 & 145,6 & -3,2\% & 303,2 & 216,4 & -28,6\% \\
2 & JavaScript & monolityczna & - & 1385 & 1096,8 & -20,8\% & 2318,8 & 1173,4 & -49,4\% \\
3 & Java & monolityczna & - & 3390 & 3630,6 & 7,1\% & 6344,4 & 7275 & 14,7\% \\
4 & Python & mikroserwisowa & gRPC & 58,6 & 55,8 & -4,8\% & 79,6 & 63,6 & -20,1\% \\
5 & JavaScript & mikroserwisowa & gRPC & 920 & 699,2 & -24,0\% & 1925,6 & 946,4 & -50,9\% \\
6 & Java & mikroserwisowa & gRPC & 2292,4 & 2523,6 & 10,1\% & 5007,8 & 4835,8 & -3,4\% \\
\hline
\end{tabular}%
}
\end{table}

\begin{figure}[htbp]
\centering
\includegraphics[width=0.82\textwidth]{images/scaling/scaling-comparison/python_scaling_comparison}
\caption{Porównanie skalowania wertykalnego i horyzontalnego dla implementacji w języku Python.}
\label{fig:compare-python}
\end{figure}

\begin{figure}[htbp]
\centering
\includegraphics[width=0.82\textwidth]{images/scaling/scaling-comparison/javascript_scaling_comparison}
\caption{Porównanie skalowania wertykalnego i horyzontalnego dla implementacji w języku JavaScript.}
\label{fig:compare-javascript}
\end{figure}

\begin{figure}[htbp]
\centering
\includegraphics[width=0.82\textwidth]{images/scaling/scaling-comparison/java_scaling_comparison}
\caption{Porównanie skalowania wertykalnego i horyzontalnego dla implementacji w języku Java.}
\label{fig:compare-java}
\end{figure}

Porównanie obu mechanizmów skalowania prowadzi do jednoznacznego wniosku, że nie istnieje
jedno uniwersalne rozwiązanie najlepsze dla wszystkich badanych technologii. W przypadku Pythona
i JavaScriptu w każdym analizowanym wariancie lepsze rezultaty dawało skalowanie horyzontalne.
Najbardziej wyraźnie widać to dla JavaScriptu, gdzie dla monolitu przewaga skalowania
horyzontalnego nad wertykalnym przy czterokrotnym zwiększeniu zasobów wyniosła 49,4\%,
a dla mikroserwisów aż 50,9\%. Oznacza to, że w tych implementacjach mnożenie instancji
przynosi znacznie lepszy efekt niż wzmacnianie pojedynczego serwera.

Java zachowywała się inaczej. Dla monolitu zarówno przy wzroście zasobów 2$\times$, jak i 4$\times$
lepsze wyniki uzyskało skalowanie wertykalne. Przewaga wyniosła odpowiednio 7,1\% oraz 14,7\%.
W przypadku mikroserwisów java sytuacja była bardziej wyrównana: dla konfiguracji 2$\times$
lepsze okazało się skalowanie wertykalne, natomiast dla konfiguracji 4$\times$ nieznacznie wyższy
wynik dało skalowanie horyzontalne. Różnica była jednak niewielka i wyniosła jedynie 3,4\%
na korzyść skalowania horyzontalnego.

Uzyskane rezultaty sugerują, że dla środowiska java zwiększanie mocy pojedynczej instancji może
być bardzo skuteczną strategią, szczególnie dla aplikacji monolitycznych. Z kolei dla Pythona i
JavaScriptu bardziej racjonalne wydaje się zwiększanie liczby instancji, ponieważ taka strategia lepiej
wykorzystuje charakterystykę tych środowisk uruchomieniowych. W praktyce oznacza to, że decyzja
o wyborze sposobu skalowania nie powinna wynikać wyłącznie z architektury systemu, lecz również
z właściwości technologii implementacyjnej.

\section{Podsumowanie wyników}

Na podstawie przeprowadzonej analizy można sformułować kilka kluczowych wniosków. Po pierwsze,
najwyższą bezwzględną wydajność w badanym zestawie osiągnęły implementacje w języku Java.
Dotyczy to zarówno architektury monolitycznej, jak i mikroserwisowej z \texttt{gRPC}. Po drugie,
architektura monolityczna we wszystkich językach osiągała wyższą maksymalną przepustowość niż
odpowiadająca jej wersja mikroserwisowa, co potwierdza wpływ narzutu komunikacyjnego na wynik
końcowy.

Po trzecie, skuteczność skalowania zależała silnie od języka programowania. Java najlepiej reagowała
na skalowanie wertykalne, natomiast Python i JavaScript wyraźnie korzystały bardziej ze skalowania
horyzontalnego. Szczególnie istotne jest to w przypadku JavaScriptu, dla którego dokładanie CPU
i pamięci do pojedynczej instancji dawało mały efekt, natomiast zwiększenie liczby instancji istotnie
podnosiło przepustowość.

Ostatecznie wyniki wskazują, że dobór strategii skalowania powinien być podejmowany na podstawie
empirycznych pomiarów dla konkretnej technologii i sposobu implementacji. Sama decyzja o wyborze
architektury monolitycznej lub mikroserwisowej nie jest wystarczająca. Równie ważne jest to, czy
wybrane środowisko uruchomieniowe lepiej wykorzystuje większą liczbę zasobów pojedynczej maszyny,
czy też większą liczbę równoległych instancji.
